\documentclass[aspectratio=169,UTF8,zihao=-4]{ctexbeamer}

\usetheme{Madrid}
\usecolortheme{default}
\setbeamertemplate{navigation symbols}{}
\setbeamertemplate{headline}{}
\setbeamertemplate{frametitle}{}
\setbeamertemplate{footline}[frame number]

\usepackage{amsmath}
\usepackage{booktabs}
\usepackage{array}

\title{范文1第1次提交\\Introduction 阅读报告}
\author{王鲲淏-22307100011}
\date{Social Research}

\newcommand{\bigidea}[1]{
  \begin{center}
    \vspace{0.8em}
    {\LARGE \bfseries #1}
    \vspace{0.6em}
  \end{center}
}

\newcommand{\sectionpageframe}[1]{
  \begin{frame}[plain]
    \begin{center}
      \vfill
      {\Huge \bfseries #1}
      \vfill
    \end{center}
  \end{frame}
}

\begin{document}

\begin{frame}[plain]
  \titlepage
\end{frame}

\sectionpageframe{论文与任务}

\begin{frame}
  \bigidea{选定论文}
  \begin{itemize}
    \item \Large Wang, X., Hong, Z., Xu, Y., Zhang, C., \& Ling, H. (2014)
    \item \Large \textit{Relevance Judgments of Mobile Commercial Information}
    \item \Large 期刊：\textit{Journal of the Association for Information Science and Technology}
    \item \Large 本周任务：阅读 Introduction，回答研究动机与研究问题
  \end{itemize}
\end{frame}

\sectionpageframe{研究动机}

\begin{frame}
  \bigidea{实践动机}
  \begin{itemize}
    \item \Large 移动商务场景下，用户会持续收到大量商业信息与移动广告。
    \item \Large 平台和商家需要判断：什么样的信息会被用户认为“相关”，什么样的信息会被认为“打扰”。
    \item \Large 如果不能提高信息相关性，就会降低点击、接受与转化效果。
    \item \Large 如果过度依赖位置与时间等个性化线索，又可能触发用户隐私担忧。
  \end{itemize}
\end{frame}

\begin{frame}
  \bigidea{理论动机}
  \begin{itemize}
    \item \Large 既有相关性判断研究主要解释文档或网页为何相关，但较少处理移动商业信息。
    \item \Large 既有移动营销研究强调位置、时间和个性化，却没有形成完整的 relevance model。
    \item \Large 因此本文试图把 relevance judgment、mobile marketing 和 privacy 三条文献整合起来。
    \item \Large 研究还把 privacy concerns 视为 affective relevance 的一种表现，扩展了相关性理论。
  \end{itemize}
\end{frame}

\sectionpageframe{研究问题}

\begin{frame}
  \bigidea{本文的核心研究问题}
  \begin{enumerate}
    \item \Large 用户如何判断一条移动商业信息是否与自己相关？
    \item \Large 位置、时间等情境因素会不会影响这种相关性判断？
    \item \Large 信息内容与情境因素又会如何影响用户的隐私担忧？
  \end{enumerate}
\end{frame}

\begin{frame}
  \bigidea{研究问题陈述可如何展示}
  \begin{itemize}
    \item \Large 原文问题句集中在摘要和 Introduction 开头。
    \item \Large 可直接展示：How do users judge the relevance of such information? Is their relevance judgment affected by contextual factors, such as location and time? How do message content and contextual factors affect users' privacy concerns?
    \item \Large 报告时先讲现实问题，再落到三个具体问题，逻辑会更清楚。
  \end{itemize}
\end{frame}

\end{document}
